\section{Квадратичное программирование в методе опорных векторов}
\label{sec:svm}

\subsection{Постановка задачи классификации}

Дана обучающая выборка $(x_i, y_i)$, $i = \overline{1,l}$, где
$x_i \in \R^n$, $y_i \in \{-1, 1\}$. Нужно построить разделяющую гиперплоскость
\[
  w\T x + b = 0 ,
\]
где $w$ — вектор нормали, $b$ — свободный член; новая точка классифицируется
правилом $\sgn\big(w\T x + b\big)$.

Гиперплоскость разделяет выборку, если $y_i\big(w\T x_i + b\big) \ge 1$ при
всех $i$ (при $y_i = 1$ это $w\T x_i + b \ge 1$, при $y_i = -1$ —
$w\T x_i + b \le -1$).

\subsubsection*{Геометрический зазор}

Расстояние от $x_i$ до гиперплоскости равно
$\rho_i = \abs{w\T x_i + b} / \norm{w}$. При выполнении
$y_i(w\T x_i + b) \ge 1$ минимальное расстояние равно $\rho = 1/\norm{w}$, а
ширина полосы между опорными гиперплоскостями $w\T x + b = \pm 1$ равна
$2/\norm{w}$. Поэтому максимизация зазора равносильна минимизации
$\tfrac12\norm{w}^2$.

\subsubsection*{Прямая задача}

Для линейно разделимой выборки
\[
  \tfrac12 \norm{w}^2 \to \min_{w,b}
  \qquad\text{при}\qquad
  y_i\big(w\T x_i + b\big) \ge 1 , \quad i = \overline{1,l} .
\]
Если выборка не разделима линейно, вводят переменные ошибок $\xi_i \ge 0$ и
\textbf{мягкий зазор}:
\begin{equation}\label{eq:svm-primal}
  \tfrac12 \norm{w}^2 + C \sum_{i=1}^{l}\xi_i \ \to\ \min_{w,b,\xi} ,
  \qquad
  y_i\big(w\T x_i + b\big) \ge 1 - \xi_i , \quad \xi_i \ge 0 ,
\end{equation}
где $C > 0$ — параметр штрафа за ошибки.

\begin{idea}
Это задача \textbf{квадратичного программирования}: целевая функция
квадратична, ограничения линейны. Чем больше $C$, тем сильнее наказываются
нарушения зазора — при $C \to \infty$ получается жёсткий зазор, при малом $C$
модель допускает много ошибок ради широкой полосы. Именно $C$ управляет
балансом между переобучением и недообучением.
\end{idea}

\subsection{Двойственная задача}

Запишем ограничения в виде $1 - \xi_i - y_i(w\T x_i + b) \le 0$,
$-\xi_i \le 0$, введём множители $\alpha_i \ge 0$ и $\mu_i \ge 0$:
\[
  L(w,b,\xi,\alpha,\mu) = \tfrac12\norm{w}^2 + C\sum_{i=1}^{l}\xi_i
    + \sum_{i=1}^{l}\alpha_i\Big(1 - \xi_i - y_i\big(w\T x_i + b\big)\Big)
    - \sum_{i=1}^{l}\mu_i \xi_i .
\]

\begin{theorem}\label{thm:svm-w}
В седловой точке функции Лагранжа
\[
  w = \sum_{i=1}^{l}\alpha_i y_i x_i .
\]
\end{theorem}

\begin{proof}
$\nabla_w L = w - \sum_i \alpha_i y_i x_i$; условие стационарности
$\nabla_w L = 0$ даёт требуемое.
\end{proof}

Остальные условия стационарности:
\[
  \frac{\partial L}{\partial b} = -\sum_{i=1}^{l}\alpha_i y_i = 0
  \ \Longrightarrow\ \sum_{i=1}^{l}\alpha_i y_i = 0 ,
  \qquad
  \frac{\partial L}{\partial \xi_i} = C - \alpha_i - \mu_i = 0
  \ \Longrightarrow\ 0 \le \alpha_i \le C
\]
(последнее — из $\mu_i \ge 0$ и $\alpha_i \ge 0$). Подставляя
$w = \sum_i \alpha_i y_i x_i$ в $L$, после сокращений получаем
\textbf{двойственную задачу}
\begin{equation}\label{eq:svm-dual}
  W(\alpha) = \sum_{i=1}^{l}\alpha_i
    - \frac12 \sum_{i=1}^{l}\sum_{j=1}^{l}
      \alpha_i\alpha_j y_i y_j\, x_i\T x_j \ \to\ \max_\alpha ,
  \qquad
  \sum_{i=1}^{l}\alpha_i y_i = 0 , \quad 0 \le \alpha_i \le C .
\end{equation}

\subsubsection*{Ядровая форма}

Заменяя скалярное произведение ядром
$K(x_i,x_j) = \varphi\T(x_i)\varphi(x_j)$, получаем
\[
  \sum_{i}\alpha_i - \frac12\sum_{i}\sum_j \alpha_i\alpha_j y_iy_j K(x_i,x_j)
  \to \max_\alpha ,
  \qquad
  \sum_i \alpha_i y_i = 0 , \quad 0 \le \alpha_i \le C ,
\]
а решающая функция —
$F(x) = \sum_{i}\alpha_i y_i K(x_i, x) + b$, классификация $\sgn F(x)$.

\begin{theorem}\label{thm:svm-convex}
Двойственная задача SVM — выпуклая задача квадратичного программирования.
\end{theorem}

\begin{proof}
Запишем $W(\alpha) = \sum_i \alpha_i - \tfrac12 \alpha\T Q \alpha$, где
$Q_{ij} = y_i y_j K(x_i,x_j)$. Для произвольного $z$, обозначив
$u_i = z_i y_i$,
\[
  z\T Q z = \sum_{i}\sum_j z_i z_j y_i y_j K(x_i,x_j)
          = \sum_{i}\sum_j u_i u_j K(x_i, x_j) \ \ge\ 0 ,
\]
так как ядровая матрица Грама положительно полуопределена. Значит
$Q \succeq 0$, функция $-\tfrac12\alpha\T Q\alpha$ вогнута, линейная часть
$\sum_i\alpha_i$ и вогнута, и выпукла, поэтому $W$ вогнута. Максимизация
вогнутой функции при линейных ограничениях — выпуклая задача.
\end{proof}

\subsubsection*{Условия ККТ и опорные векторы}

\[
  y_i\big(w\T x_i + b\big) \ge 1 - \xi_i ,\quad \xi_i \ge 0 ;
  \qquad 0 \le \alpha_i \le C ;
  \qquad \sum_i \alpha_i y_i = 0 ;
\]
\[
  \alpha_i\Big(y_i\big(w\T x_i + b\big) - 1 + \xi_i\Big) = 0 ,
  \qquad
  \big(C - \alpha_i\big)\xi_i = 0 .
\]
Отсюда:
\begin{itemize}[nosep,left=0pt]
  \item $\alpha_i = 0$ — объект не влияет на гиперплоскость;
  \item $0 < \alpha_i < C$ — объект лежит точно на границе зазора,
        $y_i(w\T x_i + b) = 1$;
  \item $\alpha_i = C$ — объект нарушает зазор или классифицирован ошибочно.
\end{itemize}
Объекты с $\alpha_i > 0$ называются \textbf{опорными векторами}.

\subsection{Алгоритмы решения}

\paragraph{Метод активного набора.}
Предполагается, что часть ограничений активна ($\alpha_i = 0$ либо
$\alpha_i = C$); на каждой итерации решается задача при фиксированном активном
наборе, проверяются условия ККТ, набор корректируется. Применим, когда матрица
квадратичной формы положительно полуопределена, ограничения линейны и можно
решать соответствующие СЛАУ. Минус: при большом $l$ активный набор часто
меняется, а матрицу ядра приходится хранить целиком.

\paragraph{Метод внутренней точки.}
Ограничения $0 < \alpha_i < C$ заменяются логарифмическим барьером
(\S\ref{subsec:barrier})
\[
  -\mu\sum_{i=1}^{l}\ln\alpha_i - \mu\sum_{i=1}^{l}\ln(C - \alpha_i) ,
\]
и решается последовательность гладких задач при $\mu \to +0$. Минус: большие
системы линейных уравнений на каждом шаге и необходимость строго внутренней
стартовой точки.

\paragraph{Методы декомпозиции.}
Оптимизация ведётся не по всем $\alpha_1,\dots,\alpha_l$ сразу, а по небольшому
рабочему множеству $B_k \subset \{1,\dots,l\}$; остальные переменные
фиксируются. Это и позволяет работать при очень больших $l$.

\subsection{Алгоритм SMO}
\label{subsec:smo}

SMO — частный случай декомпозиции, где на каждой итерации оптимизируются
\emph{ровно две} переменные $\alpha_i, \alpha_j$.

\begin{idea}
Почему именно две. Из-за ограничения $\sum_s \alpha_s y_s = 0$ изменение одной
переменной почти всегда его нарушает; двух переменных достаточно, чтобы
скомпенсировать изменение. Двумерная подзадача при этом решается аналитически,
без всякого численного метода, — отсюда крайняя дешевизна итерации.
\end{idea}

\begin{theorem}\label{thm:smo-eq}
Если новые значения выбираются так, что
$y_i\alpha_i^{\text{нов}} + y_j\alpha_j^{\text{нов}}
 = y_i\alpha_i^{\text{стар}} + y_j\alpha_j^{\text{стар}}$,
то ограничение $\sum_s \alpha_s y_s = 0$ сохраняется.
\end{theorem}

\begin{proof}
Меняются только компоненты $i$ и $j$, поэтому
\[
  \sum_s \alpha_s^{\text{нов}} y_s - \sum_s \alpha_s^{\text{стар}} y_s
  = y_i\big(\alpha_i^{\text{нов}} - \alpha_i^{\text{стар}}\big)
  + y_j\big(\alpha_j^{\text{нов}} - \alpha_j^{\text{стар}}\big) = 0
\]
по условию. Значит новая сумма равна старой, то есть нулю.
\end{proof}

\subsubsection*{Формулы SMO}

Обозначим $F(x) = \sum_s \alpha_s y_s K(x_s,x) + b$ и ошибку на объекте
$E_i = F(x_i) - y_i$. Для выбранной пары $i, j$:
\[
  \eta = K(x_i,x_i) + K(x_j,x_j) - 2K(x_i,x_j) .
\]
При $\eta > 0$ новое значение до отсечения
\[
  \alpha_j^{\text{нов}} = \alpha_j^{\text{стар}}
    + \frac{y_j\big(E_i - E_j\big)}{\eta} ,
\]
затем оно отсекается отрезком $L \le \alpha_j^{\text{нов}} \le H$:
\[
  y_i \ne y_j : \quad
  L = \max\big(0,\ \alpha_j^{\text{стар}} - \alpha_i^{\text{стар}}\big) ,
  \quad
  H = \min\big(C,\ C + \alpha_j^{\text{стар}} - \alpha_i^{\text{стар}}\big) ;
\]
\[
  y_i = y_j : \quad
  L = \max\big(0,\ \alpha_i^{\text{стар}} + \alpha_j^{\text{стар}} - C\big) ,
  \quad
  H = \min\big(C,\ \alpha_i^{\text{стар}} + \alpha_j^{\text{стар}}\big) .
\]
После этого
\[
  \alpha_i^{\text{нов}} = \alpha_i^{\text{стар}}
    + y_i y_j\big(\alpha_j^{\text{стар}} - \alpha_j^{\text{нов}}\big) .
\]

\subsubsection*{Обновление свободного члена}

\[
  b_1 = b^{\text{стар}} - E_i
    - y_i\Delta\alpha_i K(x_i,x_i) - y_j\Delta\alpha_j K(x_i,x_j) ,
\]
\[
  b_2 = b^{\text{стар}} - E_j
    - y_i\Delta\alpha_i K(x_i,x_j) - y_j\Delta\alpha_j K(x_j,x_j) .
\]
Если $0 < \alpha_i^{\text{нов}} < C$, берут $b^{\text{нов}} = b_1$; если
$0 < \alpha_j^{\text{нов}} < C$ — берут $b_2$; если обе переменные на границах
— $b^{\text{нов}} = (b_1 + b_2)/2$.

\begin{example}
Линейный SVM, $C = 1$, ядро $K(x_i,x_j) = x_i x_j$, выборка в $\R$:
\[
  x_1 = -1,\ y_1 = -1 ; \quad x_2 = 0,\ y_2 = -1 ; \quad
  x_3 = 1,\ y_3 = 1 ; \quad x_4 = 2,\ y_4 = 1 .
\]
Старт $\alpha^0 = 0$, $b^0 = 0$, тогда $F \equiv 0$ и
$E = (1;\,1;\,-1;\,-1)$.

\textbf{Итерация 1}, пара $i=1$, $j=3$: $K_{11} = K_{33} = 1$,
$K_{13} = -1$, $\eta = 4$. Так как $y_1 \ne y_3$, то $L = 0$, $H = 1$.
Получаем $\alpha_3^{\text{нов}} = 0 + \tfrac{1\cdot(1-(-1))}{4} = \tfrac12$
(в отрезке), $\alpha_1^{\text{нов}} = \tfrac12$. Проверка:
$\sum\alpha_i y_i = 0$. Обе формулы дают $b_1 = b_2 = 0$, значит $b^1 = 0$,
$w = \tfrac12 + \tfrac12 = 1$, $F(x) = x$.

\textbf{Итерация 2}, пара $i=2$, $j=3$: $E_2 = 1$, $E_3 = 0$,
$K_{22} = 0$, $K_{33} = 1$, $K_{23} = 0$, $\eta = 1$. При $y_2 \ne y_3$
получаем $L = \tfrac12$, $H = 1$, и
$\alpha_3^{\text{нов}} = \tfrac12 + 1 = \tfrac32 \to$ после отсечения
$\alpha_3^{\text{нов}} = 1$; тогда $\alpha_2^{\text{нов}} = \tfrac12$.
Итог $\alpha^2 = \big(\tfrac12;\,\tfrac12;\,1;\,0\big)$,
$\sum\alpha_iy_i = 0$, $b^2 = -1$, $w = \tfrac32$,
$F(x) = \tfrac32 x - 1$.
\end{example}

\begin{idea}
Ограничения алгоритмов: матрица ядра имеет размер $l \times l$ и при больших
$l$ не помещается в память; активные методы и методы внутренней точки для
таких выборок слишком дороги; SMO делает очень дешёвые итерации, но их может
потребоваться много, и качество сильно зависит от эвристики выбора пары
$(i,j)$ — обычно берут пару, наиболее нарушающую условия ККТ.
\end{idea}
